Intrinsic image decomposition (IID) tries to separate a photograph into two layers:
the \emph{albedo}, which is the true colour of each surface, and the \emph{shading},
which captures how light falls on those surfaces.
Most IID models assume that light is white or achromatic.
In real indoor scenes, however, light sources have colour: warm lamps cast orange tints,
while windows let in cool blue daylight. This coloured illumination
becomes entangled with the recovered albedo.
As a result, the same painted wall looks like it has different colours depending on
which light is on, even though the paint itself has not changed.

In this work we tackle this coloured illumination problem directly.
We propose \emph{Cross-Render Albedo Invariance} (CARI), a training strategy
that uses pairs of images of the same scene captured under different illuminants.
The core idea is simple: if the same surface appears in two photographs lit differently,
its albedo must be identical in both.
We enforce this with a loss that penalises any change in the predicted albedo
across the two views, while requiring that the shading layer absorbs the difference.

We implement CARI inside a model built on a frozen DINOv2-Large encoder
and a DPT-style decoder.
A key practical insight is that standard white-balance pre-processing, commonly
applied to training data, removes the very illuminant colour information that
CARI needs to function; we therefore train on raw, unprocessed image pairs.

Evaluating this claim turned out to require correcting the way it is measured. The
chroma-cast score used throughout this literature, and in earlier drafts of this work,
pools albedo chromaticity across all the materials of a scene before taking its
variance, and therefore conflates the illuminant-induced drift it intends to measure with
the ordinary chromatic diversity of the scene. Because that second term is computed from
the \emph{prediction}, a model can lower the score simply by washing the colour out of its
albedo. The metric pays models to destroy the very quantity the task exists to recover. We
decompose it into a scale-normalised invariance term and a fidelity term anchored on
measured albedo, and show that the correction reverses the ranking of our own ablation.

We evaluate on four benchmarks spanning real and synthetic scenes: MID, IIW, MAW and ARAP.
Under the corrected metric our model is the only one in the comparison that is
simultaneously colour-faithful and competitively invariant: it reproduces the
between-material chroma of measured ground-truth albedo almost exactly (a fidelity of
$1.008$, against $0.825$ for the strongest diffusion baseline and $0.575$ for the strongest
feed-forward one) and attains the best material-lightness stability in the roster, while
remaining within a fifth of the most illuminant-invariant method on hue. We do not claim
the invariance crown: two external methods hold it, and we show that both purchase it by
discarding roughly a fifth and two fifths of real albedo colour respectively. The same
degenerate trade that our own colour-skip ablation makes, which is visible to the naked eye.
These results are obtained with $18.5$M trained parameters and a single forward pass,
some seventy times fewer trained weights and five times faster than the diffusion baselines.
The remaining gap is in perceptual shading quality measured by IIW, which we analyse through
the constancy--structure tension inherent to IID.
